\chapter{Code Snippets} \label{ap1:loren}
 
\section{Pulse Generation}

\begin{lstlisting}[language=C, caption={Buffer processing function.}, label={lst:processBuffer}]
    void process_buffer(float32_t *mainBuf) {
        int16_t *rx_buf, *tx_buf;
    
        if (rx_buffer_proc == PING)
            rx_buf = (int16_t *)PING_IN;
        else
            rx_buf = (int16_t *)PONG_IN;
    
        if (tx_buffer_proc == PING)
            tx_buf = (int16_t *)PING_OUT;
        else
            tx_buf = (int16_t *)PONG_OUT;
    
        int16_t sampleL, sampleR;
        for (int i = 0; i < (PING_PONG_BUFFER_SIZE); i++){
            // read both input samples
            sampleL = *(rx_buf++); sampleR = *(rx_buf++);
    
            // copy the input buffer to the output buffer
            *(tx_buf++) = (int16_t) floor((float32_t) 0.5 + mainBuf[i]);
            *(tx_buf++) = sampleR; // RIGHT channel remains untouched
        }
    
        RX_buffer_full = 0;
        TX_buffer_empty = 0;
    }
\end{lstlisting}

\begin{lstlisting}[language=C, caption={Pulse generation \texttt{main()} funciton.}, label={lst:mainPulseGen}]
int main(void) {
    uint16_t idxMainBuf = 0, vibrato = 0, contourIndex = 0; // reset main buffer index
    uint16_t gPulseLen, lastPulseLen, filteredLength;

    float32_t mainBuf[2 * PING_PONG_BUFFER_SIZE] = {0.0f};
    float32_t gPulseBuf[MAX_PULSE_BUFFER] = {0.0f}; // 220 ~ 22050/100

    stm32f7_wm8994_init(AUDIO_FREQUENCY_22K, // 22050 Hz
                        IO_METHOD_DMA,
                        INPUT_DEVICE_INPUT_LINE_1,
                        OUTPUT_DEVICE_HEADPHONE,
                        WM8994_HP_OUT_ANALOG_GAIN_0DB,
                        WM8994_LINE_IN_GAIN_0DB,
                        WM8994_DMIC_GAIN_9DB,
                        SOURCE_FILE_NAME,
                        NOGRAPH);

    // generate single prototype pulse 10% shorter than that of the fundamental period
    gPulseLen = genPulse(FS, F0 * 1.1, gPulseBuf, PULSE_GAIN);

    for(;;){
        while (!(RX_buffer_full && TX_buffer_empty)){}

        // 1. Copy first half to the output during the DMA processing
        process_buffer(mainBuf); 

        // 2. Copy second half to first half
        for (size_t i = 0; i < N; i++){
            mainBuf[i] = mainBuf[i + N];
            mainBuf[i + N] = 0.0f;
        }

        // Place pulses inside the buffer
        while (idxMainBuf < N){
            lastPulseLen = getPulseLen(FS, contour[contourIndex]);

            // add vibrato
            lastPulseLen += (uint16_t) floor((float32_t) 0.5f + (float32_t) 7.0f * sin((float32_t)2.0f * PI * vibrato / (float32_t)28.0f));

            // 3. Filter and accumulate the pulse into the mainbuffer
            filteredLength = lastPulseLen + 20; // dummy

            // 4. Copy *filtered* pulse to main buffer
            for (size_t i = 0; i < filteredLength; i++)
                mainBuf[idxMainBuf++] += gPulseBuf[i];

            idxMainBuf -= (filteredLength - lastPulseLen);
            
            vibrato = (vibrato+1) % 28;
            contourIndex = (contourIndex < contourLen) ? contourIndex + 1 : contourIndex;
        }

        idxMainBuf -= N;
    }
}
\end{lstlisting}
    

\section{Filtering}

\begin{lstlisting}[language=C, caption={Vowel Synthesis in C}, label={lst:vowelSynth}]
int main(void) {
    // initialize the board and the audio codec
    stm32f7_wm8994_init(AUDIO_FREQUENCY_22K, // 22050 Hz
                    IO_METHOD_DMA,
                    INPUT_DEVICE_INPUT_LINE_1,
                    OUTPUT_DEVICE_HEADPHONE,
                    WM8994_HP_OUT_ANALOG_GAIN_0DB,
                    WM8994_LINE_IN_GAIN_0DB,
                    WM8994_DMIC_GAIN_9DB,
                    SOURCE_FILE_NAME,
                    NOGRAPH);

    uint16_t idxMainBuf = 0, vibrato = 0, contourIndex = 0; // reset main buffer index
    uint16_t gPulseLen, lastPulseLen, filteredLength;

    float32_t mainBuf[2 * PING_PONG_BUFFER_SIZE] = {0.0f};
    float32_t gPulseBuf[MAX_PULSE_BUFFER] = {0.0f}; // 220 ~ 22050/100
    float32_t colorPulse[MAX_PULSE_BUFFER] = {0.0f};
        
    float32_t tilt[N] = {0.0f};

    // * generate single prototype pulse 10% shorter than that of the fundamental period
    //gPulseLen = genPulse(FS, F0 * 1.1f, gPulseBuf, PULSE_GAIN);
    gPulseLen = genPulse(FS, F0, gPulseBuf, PULSE_GAIN);
        
    // pre compute the pulse spectral tilt of the filter (1/lfgain)
    tiltGain(tilt, N, FS, F0);
    
        // * precompute sin window and odft exponentials
    for (uint32_t i = 0; i < N; i++) {
        sinW[i] = (float32_t) sin(PI / N * (i + 0.5f));
            
        dirExp[i].real = (float32_t) cos(-PI * i / N);
        dirExp[i].imag = (float32_t) sin(-PI * i / N);

        invExp[i].real = (float32_t) cos(PI * i / N);
        invExp[i].imag = (float32_t) sin(PI * i / N);
    }

    float32_t autocorrelation[LPCorder + 1] = {0.0f};
    float32_t aCoeff[LPCorder + 1] = {0.0f};
    float32_t error = 0.0f, b0 = 0.0f;

    float32_t psd[N] = {0.0f};
    
    for(;;){
        while (!(RX_buffer_full && TX_buffer_empty)){}

        // 1. Copy first half to the output during the DMA processing
        process_buffer(mainBuf);

        arm_mult_f32(psd, tilt, psd, N); // Multiply the PSD by the tilt gain factor	
                
        for (uint32_t i = 0; i < N; i++){ // create complex array from PSD
            tempBuff[i].real = psd[i];
            tempBuff[i].imag = 0.0f;
        }
                
        // Use weiner khinchine (apply IODFT) to compute the autocorrelation sequence from the PSD
        arm_cfft_f32(&arm_cfft_sR_f32_len1024, (float32_t *)(tempBuff), 1, 1); // ifft flag = 1, doBitReverse = 1
        arm_cmplx_mult_cmplx_f32((float32_t *)(tempBuff), (float32_t *)(invExp), (float32_t *)(currBuff), N); // iodft

        // only first (LPCorder + 1) values are needed for levinson recursion
        for (uint32_t i = 0; i <= LPCorder; i++)
            autocorrelation[i] = tempBuff[i].real;    // assuming here the complex part will be zero

        levinson(autocorrelation, LPCorder, aCoeff, &error); // LPC analysis

        arm_sqrt_f32(error, &b0); // compute b0 gain factor

        // 2. Copy second half to first half
        for (int32_t i = 0; i < N2; i++){
            mainBuf[i] = mainBuf[i + N2];
            mainBuf[i + N2] = 0.0f;
        }

        // Place pulses inside the buffer
        while (idxMainBuf < N2){
            // follow f0 contourIndex
            lastPulseLen = getPulseLen(FS, F0); // dummy pulse length just for placement

            // 3. Filter and accumulate the pulse into the mainbuffer
            filteredLength = MAX_PULSE_BUFFER;
            iirDirectTypeII(gPulseBuf, colorPulse, MAX_PULSE_BUFFER, b0, aCoeff, LPCorder + 1);

            // 4. Copy *filtered* pulse to main buffer
            for (int32_t i = 0; i < filteredLength; i++)
                mainBuf[idxMainBuf++] += 0.00002f * colorPulse[i];

            idxMainBuf -= (filteredLength - lastPulseLen);
        }

        idxMainBuf -= N2;
    }
}
\end{lstlisting}

\section{Overlap Add conversion}

\begin{lstlisting}[language=Python, caption={Complete System evaluation in python}, label={lst:pythonConversion}, style=mystyle]
    N = 1024; N2 = N // 2; INT16_MAX = 32767
    
    # hanning sine window for each frame
    sinW = np.sin(np.pi * (np.arange(0, N) + 0.5) / N)
    
    # input and output buffers
    old_inBuff  = np.zeros(N2)
    old_outBuff = np.zeros(N2)
    
    # pulse buffers
    pulseBuffer = np.zeros(N) # 1024 samples long
    # initialization for the first the pulse, such that the alpha coefficient for the PSD calculation is never 0
    lastPulsePos = 1 
    
    # power spectral density
    old_psd = np.zeros(N)
    pulse_psd = np.zeros(N)
    
    # output file
    outputFile = np.array([])
    # tilt compensation factor
    tilt = tiltGain(F0, FS, N)
    tilt = np.concatenate((tilt, tilt[::-1]))
    
    aaa = np.array([])
    
    currF0 = 0
    f0contour = np.concatenate((contour, np.full(10*FS, contour[-1])))
    
    # load new frames 512 samples at a time
    for start in range(0, len(datain) - N2, N2):
        in_frame = datain[start : start + N2]
    
        if(len(in_frame) != N2):
            in_frame = np.pad(in_frame, (0, N2 - len(in_frame)), 'constant')
    
        # Concatenate input sequence and apply window
        winFrame = np.concatenate((old_inBuff, in_frame)) * sinW
        old_inBuff = in_frame
    
        # DFT and compute PSD for current frame
        curr_spectrum = odft(winFrame, N)
        curr_psd = np.abs(curr_spectrum) ** 2
    
        # pulse generation
        pulseBuffer[:N2] = pulseBuffer[N2:] # shift the buffer
        pulseBuffer[N2:] = 0 # clear the second half of the buffer
    
        # compute centre of gravity and energy from PSD
        cog = np.sum(curr_psd[:N2] * (np.arange(0, N2)*FS)/N) / np.sum(curr_psd[:N2])
        energy = np.sum(curr_psd[:N2])
        
        unvoiced = m(cog, Fvuv, Fuv)
        activity = WAD(in_frame) # here we need the in_frame such that we can compute the zero crossing rate
        
        aaa = np.concatenate((aaa, np.array([unvoiced])))
    
        while lastPulsePos < N2:
            # generate both a pulse and noise
            f0 = f0contour[currF0]
            pLen = int(FS/f0)
            t = np.arange(0, 1/f0, 1/FS)
            idx = np.arange(currF0, currF0 + len(t))
            pulse = v_glotlf(1, t*f0contour[idx])
            noise = np.random.normal(0, 0.1, len(pulse)) * 0.2
    
            alpha = lastPulsePos / N2# psd mixing coefficient
            pulse_psd = (alpha)*curr_psd + (1-alpha)*old_psd # average the psds
            pulse_psd = pulse_psd / tilt # tilt compensation
            
            uvmix = unvoiced * 0.5 + 0.5
            pulse = (pulse * (1-uvmix) + noise*uvmix) * activity # add noise if unvoiced
            
            r = iodft(pulse_psd, N).real # weiner-khinchin
            [a, e, k] = levinson(r, 22)
            b0 = np.sqrt(e)
    
            lastPulseLen = len(pulse)
            pulse = np.concatenate((pulse, np.zeros(300)))
    
            # filter the pulse
            filteredPulse = signal.lfilter(b0, a, pulse)
            pulseBuffer[lastPulsePos : lastPulsePos + len(filteredPulse)] += filteredPulse
    
            # update pulse position
            lastPulsePos += lastPulseLen
            currF0 += pLen # update curr frequency
    
        if(unvoiced > 0.5 or activity < 0.5): # Reset frequency contour if unvoiced
            currF0 = 0
    
        lastPulsePos -= N2
        output = pulseBuffer[:N2] # first half of the buffer
        old_psd = curr_psd
    
        # Add the output through overlap add and save the previous output
        outputFile = np.concatenate((outputFile, output))
        old_outBuff = output
    
    outputFile = outputFile / max(abs(outputFile)) * INT16_MAX
\end{lstlisting}

\begin{lstlisting}[language=C, caption={Speech conversion with segmentation in C}, label={lst:segmentationC}]
    int main(void) {
        // initialize the board and the audio codec
        stm32f7_wm8994_init(AUDIO_FREQUENCY_22K, // 22050 Hz
                        IO_METHOD_DMA,
                        INPUT_DEVICE_INPUT_LINE_1,
                        OUTPUT_DEVICE_HEADPHONE,
                        WM8994_HP_OUT_ANALOG_GAIN_0DB,
                        WM8994_LINE_IN_GAIN_0DB,
                        WM8994_DMIC_GAIN_9DB,
                        SOURCE_FILE_NAME,
                        NOGRAPH);
    
        uint16_t idxMainBuf = 0, vibrato = 0, contourIndex = 0; // reset main buffer index
        uint16_t gPulseLen, lastPulseLen, filteredLength;
    
        float32_t mainBuf[2 * PING_PONG_BUFFER_SIZE] = {0.0f};
        float32_t gPulseBuf[MAX_PULSE_BUFFER] = {0.0f}; // 220 ~ 22050/100
        float32_t gNoiseBuf[MAX_PULSE_BUFFER] = {0.0f};
        float32_t colorPulse[MAX_PULSE_BUFFER] = {0.0f};
            
        float32_t tilt[N] = {0.0f};
    
        // * generate single prototype pulse 10% shorter than that of the fundamental period
        gPulseLen = genPulse(FS, F0, gPulseBuf, PULSE_GAIN);
            
        // pre compute the pulse spectral tilt of the filter (1/lfgain)
        tiltGain(tilt, N, FS, F0);
        
        // * precompute sin window and odft exponentials
        for (uint32_t i = 0; i < N; i++) {
            sinW[i] = (float32_t) sin(PI / N * (i + 0.5f));
                
            dirExp[i].real = (float32_t) cos(-PI * i / N);
            dirExp[i].imag = (float32_t) sin(-PI * i / N);
    
            invExp[i].real = (float32_t) cos(PI * i / N);
            invExp[i].imag = (float32_t) sin(PI * i / N);
        }
    
        float32_t autocorrelation[LPCorder + 1] = {0.0f};
        float32_t aCoeff[LPCorder + 1] = {0.0f};
        float32_t error = 0.0f, b0 = 0.0f;
    
        float32_t currPSD[N] = {0.0f};
        float32_t lastPSD[N] = {0.0f};
        float32_t pAlpha = 0.0f;
        
        for(;;){
            while (!(RX_buffer_full && TX_buffer_empty)){}
    
            //* 1. Copy first half to the output during the DMA processing
            process_buffer(mainBuf);
    
            //* 1.5. Compute filter from currBuff
            arm_cmplx_mult_cmplx_f32((float32_t *)(currBuff), (float32_t *)(dirExp), (float32_t *)(tempBuff), N); //odft exponential //working fine
            arm_cfft_f32(&arm_cfft_sR_f32_len1024, (float32_t *)(tempBuff), 0, 1); // ifft flag = 0, doBitReverse = 1
            
            arm_cmplx_mag_squared_f32((float32_t *)(tempBuff), (float32_t *)(currPSD), N); // Compute PSD 
    
            //* 2. Copy second half to first half
            for (int32_t i = 0; i < N2; i++){
                mainBuf[i] = mainBuf[i + N2];
                mainBuf[i + N2] = 0.0f;
            }
    
            // compute the PSD's centre of gravity
            float32_t cog = 0.0f, voiced = 0.0f;
            cog = centreOfGravity(currPSD, N, FS);
            voiced = mapping(cog, 2500, 4200, 0.99);
    
            // Place pulses inside the buffer
            while (idxMainBuf < N2){
                lastPulseLen = getPulseLen(FS, F0); // dummy pulse length just for placement
                pAlpha = (idxMainBuf != 0) ? idxMainBuf/N2 : 1/N2 ; // pulse filtering coefficient (check for zero)
    
                //* 2.5. Generate a new filter based on the current pulse position
                for (uint32_t i = 0; i < N; i++){
                    tempBuff[i].real = (pAlpha * currPSD[i] + (1 - pAlpha) * lastPSD[i]) * tilt[i]; // apply tilt gain
                    tempBuff[i].imag = 0.0f;
    
                    lastPSD[i] = currPSD[i]; // save last psd
                }
    
                // Use weiner khinchine (apply IODFT) to compute the autocorrelation sequence from the PSD
                arm_cfft_f32(&arm_cfft_sR_f32_len1024, (float32_t *)(tempBuff), 1, 1); // ifft flag = 1, doBitReverse = 1
                arm_cmplx_mult_cmplx_f32((float32_t *)(tempBuff), (float32_t *)(invExp), (float32_t *)(currBuff), N); // iodft
    
                // only first (LPCorder + 1) values are needed for levinson recursion
                for (uint32_t i = 0; i <= LPCorder; i++)
                    autocorrelation[i] = tempBuff[i].real;    // assuming here the complex part will be zero
    
                levinson(autocorrelation, LPCorder, aCoeff, &error); // LPC analysis
                b0 = sqrtf(error); // gain factor
    
                //* Add noise to the pulse before filtering
                float32_t uvMix = 0.5f + voiced * 0.5f;
    
                for (uint32_t i = 0; i < MAX_PULSE_BUFFER; i++)
                    gNoiseBuf[i] = uvMix * gPulseBuf[i] + (1.0f - uvMix) * randFloat(); // add noise to the pulse
     
                //* 3. Filter and accumulate the pulse into the mainbuffer
                filteredLength = MAX_PULSE_BUFFER;
                iirDirectTypeII(gPulseBuf, colorPulse, MAX_PULSE_BUFFER, b0, aCoeff, LPCorder + 1);
    
                //* 4. Copy filtered pulse to main buffer
                for (int32_t i = 0; i < filteredLength; i++)
                    mainBuf[idxMainBuf++] += 0.00002f * colorPulse[i];
    
                // update main buffer index
                idxMainBuf -= (filteredLength - lastPulseLen);
            }
    
            for (uint32_t i = 0; i < N; i++)
                lastPSD[i] = currPSD[i];
            
            idxMainBuf -= N2;
        }
    }
\end{lstlisting}

\begin{lstlisting}[language=C, caption={Whispered activity detection algorithm in C}, label={lst:wadC}]
/// @brief Computes the Weighted Average Decision (WAD) based on energy and zero-crossing rate.
/// @param frame Pointer to an array of 512 float32_t elements representing the signal frame.
/// @param z1 Energy threshold below which the function returns 0.0f.
/// @param z2 Energy threshold above which the function returns 1.0f.
/// @param zCross Zero-crossing rate threshold to decide the WAD when energy is between z1 and z2.
/// @return Returns 1.0f if the energy is above z2 or if the zero-crossing rate is above zCross when energy is between z1 and z2. Returns 0.0f otherwise.
float32_t WAD(float32_t* frame, float32_t z1, float32_t z2, float32_t zCross) {
    float energy = 0.0f;
    int zeroCrossings = 0;
    
    for (int i = 0; i < 512; i++)
        energy += frame[i] * frame[i];
    
    // Check energy against thresholds
    if (energy >= z2)
        return 1.0f;
    else if (energy < z1)
        return 0.0f;

    // Calculate the zero-crossing rate
    for (int i = 0; i < 511; i++)
        if (frame[i] * frame[i + 1] < 0)
            zeroCrossings++;
    
    float zeroCrossingRate = (float)zeroCrossings / 512.0f;
    
    // Return WAD based on zero-crossing rate
    return (zeroCrossingRate > zCross) ? 1.0f : 0.0f;
}
\end{lstlisting}